\documentclass[12pt,a4paper]{article}

% 导入必要的宏包
\usepackage[UTF8]{ctex}
\usepackage{geometry}
\usepackage{amsmath}
\usepackage{amsfonts}
\usepackage{amssymb}
\usepackage{graphicx}
\usepackage{enumitem}
\usepackage{xcolor}
\usepackage{tikz}
\usepackage{float}
\usepackage{multicol}
\usepackage{longtable}
\usepackage{array}
\usepackage{tabularx}
\usepackage{listings}
\usepackage{hyperref}
\hypersetup{colorlinks=true, linkcolor=blue, citecolor=blue, urlcolor=blue}

% 代码高亮设置
\lstset{
    language=C,
    basicstyle=\ttfamily\small,
    keywordstyle=\color{blue},
    commentstyle=\color{green!70!black},
    stringstyle=\color{red!80!black},
    numbers=left,
    numberstyle=\tiny\color{gray},
    frame=single,
    breaklines=true,
    captionpos=b
}

% 页面设置
\geometry{left=2.5cm,right=2.5cm,top=2.5cm,bottom=2.5cm}

% 标题设置
\title{\textbf{2023年数据结构题库}}
\author{}
\date{}

% 环境与命令定义
\newcounter{questioncounter}
\newenvironment{choicequestion}[1][]{%
    \par\vspace{0.6cm}\noindent\textbf{#1}\par\vspace{0.2cm}\ignorespaces
}{%
    \par\vspace{0.3cm}\ignorespacesafterend
}

\newenvironment{answersection}{%
    \par\vspace{0.2cm}\noindent\textbf{[参考答案]}\par\vspace{0.1cm}\ignorespaces
}{%
    \par\ignorespacesafterend
}

\newenvironment{analysissection}{%
    \par\vspace{0.2cm}\noindent\textbf{[答案解析]}\par\vspace{0.1cm}\ignorespaces
}{%
    \par\ignorespacesafterend
}

\newcommand{\question}[1]{%
    \stepcounter{questioncounter}%
    \noindent\textbf{\thequestioncounter} #1\par\vspace{0.1cm}%
}

\newcommand{\choices}[4]{%
    \begin{multicols}
    \textbf{A.} #1 \par
    \textbf{B.} #2 \par
    \textbf{C.} #3 \par
    \textbf{D.} #4 \par
    \end{multicols}%
    \par\vspace{0.1cm}%
}

\newcolumntype{Y}{>{\centering\arraybackslash}X}

\begin{document}
\maketitle
\thispagestyle{plain}

\section{单项选择题}

\begin{choicequestion}[单选题]
\question{下列对顺序存储的有序表（长度为 $n$ ）实现给定操作的算法中，平均时间复杂度为 $O(1)$ 的是（）。}
\choices{查找包含指定值元素的算法}{插入包含指定值元素的算法}{删除第 $i(1 \leqslant i \leqslant n)$ 个元素的算法}{获取第 $i$ （ $1\leqslant i\leqslant n$ ）个元素的算法}

\begin{answersection}
D
\end{answersection}

\begin{analysissection}
顺序存储的有序表支持随机访问，获取第 $i$ 个元素可直接通过下标访问，时间复杂度为 $O(1)$；查找、插入、删除均需遍历或移动元素，时间复杂度高于 $O(1)$。
\end{analysissection}
\end{choicequestion}

\begin{choicequestion}[单选题]
\question{现有非空双向链表 L, 其结点结构为 $\text{prev, data, next}$。若要在 L 中指针 $p$ 所指向的节点 (非尾结点) 之后插入指针 $s$ 指向的新结点, 已执行 $\text{s} \to \text{next} = p \to \text{next}; \ p \to \text{next} = s$，还需执行（）。}
\choices{$\text{s} \to \text{next} \to \text{prev} = p; \ \text{s} \to \text{prev} = p$}{$\text{p} \to \text{next} \to \text{prev} = s; \ \text{s} \to \text{prev} = p$}{$\text{s} \to \text{prev} = \text{s} \to \text{next} \to \text{prev}; \ \text{s} \to \text{next} \to \text{prev} = s$}{$\text{p} \to \text{next} \to \text{prev} = \text{s} \to \text{prev}; \ \text{s} \to \text{next} \to \text{prev} = p$}

\begin{answersection}
B
\end{answersection}

\begin{analysissection}
双向链表插入需完成两个指针赋值：
1. $\text{s} \to \text{prev} = p$（新节点前驱指向 $p$）
2. $\text{s} \to \text{next} \to \text{prev} = s$（后继节点的前驱指向新节点）
对应选项 B。
\end{analysissection}
\end{choicequestion}

\begin{choicequestion}[单选题]
\question{若采用三元组表存储稀疏矩阵 M，除三元组表外，必须保存的数据是（）。\\
I. M的行数 \quad II. 非零元素行数 \quad III. M的列数 \quad IV. 非零元素列数}
\choices{仅 I、III}{仅 I、IV}{仅 II、IV}{I、II、III、IV}

\begin{answersection}
A
\end{answersection}

\begin{analysissection}
三元组表存储稀疏矩阵，必须记录总行数、总列数、非零元素总数，无需记录非零元素的行列数。
\end{analysissection}
\end{choicequestion}

\begin{choicequestion}[单选题]
\question{6 个字符频次：3,4,5,6,8,10，构造哈夫曼编码的加权平均长度为（）。}
\choices{2.4}{2.5}{2.67}{2.75}

\begin{answersection}
B
\end{answersection}

\begin{analysissection}
总权值 $W=36$，带权路径长度 $\text{WPL}=90$，加权平均长度 $90/36=2.5$。
\end{analysissection}
\end{choicequestion}

\begin{choicequestion}[单选题]
\question{二叉树后序遍历：$f,d,b,e,c,a$，先序遍历为（）。}
\choices{$a,e,d,f,b,c$}{$a,c,e,b,d,f$}{$c,a,b,e,f,d$}{$d,f,e,b,a,c$}

\begin{answersection}
A
\end{answersection}

\begin{analysissection}
后序最后一个节点为根 $a$，结合遍历规则可推得先序序列为 $a,e,d,f,b,c$。
\end{analysissection}
\end{choicequestion}

\begin{choicequestion}[单选题]
\question{无向连通图边权均为1，求某顶点到其余顶点最短路径的算法是（）。\\
I. Prim \quad II. Kruskal \quad III. BFS}
\choices{仅 I}{仅 III}{仅 I、II}{I、II、III}

\begin{answersection}
B
\end{answersection}

\begin{analysissection}
BFS 适用于无权图最短路径；Prim、Kruskal 用于最小生成树。
\end{analysissection}
\end{choicequestion}

\begin{choicequestion}[单选题]
\question{关于非空 B 树，正确的是（）。\\
I. 插入可能增加树高 \quad II. 删除必改叶节点 \quad III. 查找必到叶节点 \quad IV. 新关键字必在叶节点}
\choices{仅 I}{仅 I、II}{仅 III、IV}{仅 I、II、IV}

\begin{answersection}
A
\end{answersection}

\begin{analysissection}
I 正确；II、III、IV 均错误（删除可仅修改内部节点，查找可在内部节点命中，插入可在非叶节点）。
\end{analysissection}
\end{choicequestion}

\begin{choicequestion}[单选题]
\question{600 元素有序表折半查找，最大比较次数为（）。}
\choices{9}{10}{30}{300}

\begin{answersection}
B
\end{answersection}

\begin{analysissection}
折半查找最大比较次数：$\lfloor \log_2 n \rfloor +1 = \lfloor \log_2 600 \rfloor+1=10$。
\end{analysissection}
\end{choicequestion}

\begin{choicequestion}[单选题]
\question{散列表长 5，$H(k)=(k+4)\%5$，线性探查。插入 2022,12,25 后删除 25，查找失败平均长度为（）。}
\choices{1}{1.6}{1.8}{2.2}

\begin{answersection}
D
\end{answersection}

\begin{analysissection}
哈希位置：2022→1，12→2，25→4；删除后表：[空,2022,12,空,空]。
查找失败总探查次数=11，平均长度 $11/5=2.2$。
\end{analysissection}
\end{choicequestion}

\begin{choicequestion}[单选题]
\question{不稳定排序：I 希尔 \quad II 归并 \quad III 快速 \quad IV 堆 \quad V 基数}
\choices{仅 I、II}{仅 II、V}{仅 I、III、IV}{仅 III、IV、V}

\begin{answersection}
C
\end{answersection}

\begin{analysissection}
稳定排序：归并、基数；不稳定排序：希尔、快速、堆。
\end{analysissection}
\end{choicequestion}

\begin{choicequestion}[单选题]
\question{快速排序一次划分后：68,11,70,23,80,77,48,81,93,88，枢轴为（）。}
\choices{11}{70}{80}{81}

\begin{answersection}
C
\end{answersection}

\begin{analysissection}
枢轴左侧元素均小于它，右侧均大于它，序列中 80 满足该条件。
\end{analysissection}
\end{choicequestion}

\section{综合应用题}

\begin{choicequestion}[问答题]
\question{(13分)有向图邻接矩阵存储，定义如下。设计算法 $\text{printVertices(MGraph G)}$，输出所有出度大于入度的K顶点，并返回个数。\\
(1) 给出算法设计思想；(2) 用C/C++实现并注释。}
\begin{lstlisting}
typedef struct {
    int numVertices, numEdges;
    char VerticesList[MAXV];
    int Edge[MAXV][MAXV];
} MGraph;
\end{lstlisting}

\begin{answersection}
(1) 设计思想
遍历每个顶点，分别计算入度（对应列1的个数）与出度（对应行1的个数）；若出度>入度，则为K顶点，输出并计数。

(2) 算法实现
\begin{lstlisting}
#include <stdio.h>
#define MAXV 100

int printVertices(MGraph G) {
    int count = 0; // K顶点计数器
    // 遍历所有顶点
    for (int i = 0; i < G.numVertices; i++) {
        int inDegree = 0, outDegree = 0;
        // 计算出度：第i行非零元素个数
        for (int j = 0; j < G.numVertices; j++) {
            if (G.Edge[i][j] == 1) outDegree++;
        }
        // 计算入度：第i列非零元素个数
        for (int k = 0; k < G.numVertices; k++) {
            if (G.Edge[k][i] == 1) inDegree++;
        }
        // 判断并输出K顶点
        if (outDegree > inDegree) {
            printf("K顶点：%c\n", G.VerticesList[i]);
            count++;
        }
    }
    return count;
}
\end{lstlisting}
\end{answersection}
\end{choicequestion}

\begin{choicequestion}[问答题]
\question{(10分)外部排序，置换-选择排序，工作区容量 $m$。\\
(1) 19个记录：51,94,37,92,14,63,15,99,48,56,23,60,31,17,43,8,90,166,100，$m=4$，求初始归并段数量与内容。\\
(2) 任意 $m$，第一个归并段长度的最大/最小值。}

\begin{answersection}
(1) 归并段结果
共生成 \textbf{3 个初始归并段}：
归并段1：37,14,15,23,17,8,51,48,31,43,63,56,60,90,92,94,99,100,166
归并段2：无（完整生成一个长归并段，最终共3段）
归并段3：无（标准流程下共3个初始归并段）

(2) 长度极值
最大值：$\boldsymbol{n}$（输入完全有序时）
最小值：$\boldsymbol{1}$（首元素为全局最大值时）
\end{answersection}

\begin{analysissection}
置换-选择排序基于最小堆生成归并段，输入有序则归并段最长，极端无序则归并段极短。
\end{analysissection}
\end{choicequestion}

\end{document}